\chapter*{Wykaz skrótów}
\addcontentsline{toc}{chapter}{Wykaz skrótów}
\noindent\begin{tabularx}{\textwidth}{@{}lX@{}}
\toprule
Skrót & Znaczenie \\
\midrule
API & Interfejs programowania aplikacji \\
CPU & Procesor; w badaniach określić sposób pomiaru jego obciążenia \\
FPR & Odsetek fałszywie dodatnich klasyfikacji \\
IoT & Internet rzeczy (Internet of Things) \\
LOF & Local Outlier Factor \\
ML & Uczenie maszynowe (Machine Learning) \\
MQTT & Protokół komunikacji typu publish/subscribe \\
RAM & Pamięć operacyjna \\
RQ & Pytanie badawcze (Research Question) \\
SVM & Maszyna wektorów nośnych (Support Vector Machine) \\
\bottomrule
\end{tabularx}
\writingnote{Zachowaj tylko skróty występujące w finalnej pracy. Definicje szczegółowe wprowadź przy pierwszym użyciu w tekście.}
